\documentclass[12pt]{article}
\usepackage[utf8]{inputenc}
\usepackage[T1]{fontenc}
\usepackage{hyperref}
\title{El Alto Zapotec Language Resources for K-12 Curriculum Development}
\author{Compiled by Claude Sonnet 5.5 (claude-sonnet-5-5) for Andr\'as Kornai}
\date{}
\begin{document}
\maketitle

\section{Introduction}
El Alto Zapotec (ISO 639-3 zpp; Glottocode elal1235) is an Otomanguean language of Oaxaca, Mexico. It is spoken in San Pedro el Alto, San Antonino el Alto and San Andr\'es el Alto. Published speaker estimates differ: about 900 (1990 census, via Wikipedia) and 1,400 (Joshua Project). This search found very little. No school materials, grammar, dictionary or corpus turned up.

\section{Primary Educational Resources}
No primers, textbooks or school programmes for El Alto Zapotec were found in the search results.

\section{Linguistic and Reference Materials}
\begin{itemize}
\item Wikipedia overview of the language: \url{https://en.wikipedia.org/wiki/El_Alto_Zapotec}
\item Glottolog entry, whose bibliography is the best place to look for further documentation: \url{https://glottolog.org/resource/languoid/id/elal1235}
\item Wikitongues language page, which lists the language as unwritten: \url{https://wikitongues.org/languages/zpp/}
\end{itemize}

\section{Religious and Community Resources}
\begin{itemize}
\item Joshua Project profile, which reports no Bible translation: \url{https://joshuaproject.net/index.php/languages/zpp}
\end{itemize}

\section{Digital and Audio Resources}
\begin{itemize}
\item Global Recordings Network evangelism and Bible-teaching audio, with MP3 and slideshow downloads: \url{https://www.globalrecordings.net/en/language/zpp}
\end{itemize}

\section{Curriculum Development Context}
The language is reported as unwritten and is not known to be taught in schools. A native-language curriculum would first need an orthography, community-led documentation and teaching materials. Only evangelistic audio was found, and it is not educational material.
\end{document}